% ---------------------------------------------------------------------
\clearpage
\section{Reproducibilidad}\label{ap:reproducibilidad}

Todo el trabajo se reproduce con el repositorio público \url{https://github.com/CesarCapelo/BrugadaInteligenciaDeNegocio}. El conjunto de datos se descarga automáticamente de PhysioNet. Las semillas aleatorias están fijadas (semilla 42, con un desplazamiento determinista por \emph{fold}), y la huella SHA-256 de los ficheros de señal (prefijo \var{d714832551a217e6}) garantiza que los resultados corresponden a la misma versión de los datos.

\subsection*{Entorno}

Python 3.12, NumPy 1.26.4, SciPy 1.13.1, scikit-learn 1.5.2, NeuroKit2 0.2.10, WFDB 4.3.1, umap-learn 0.5.7, pandas 2.2.3 y Matplotlib 3.9.4. La lista completa está en \var{requirements.txt}.

\subsection*{Comandos}

\begin{center}
\small
\begin{tabularx}{\linewidth}{@{}p{4.3cm}L@{}}
\toprule
\textbf{Análisis} & \textbf{Comando (desde \var{src/})} \\
\midrule
Descarga de los datos & \var{bash ../scripts/download\_data.sh} \\
ETL y control de calidad (\cref{cap:datos}) & \var{python -m brugada.experiments.run\_etl} \\
Comparación del preprocesado (\cref{cap:preprocesado}) & \var{python -m brugada.experiments.run\_preprocessing} \\
Selección de instancias (\cref{cap:instancias}) & \var{python -m brugada.experiments.run\_instance\_selection} \\
Características y su selección (\cref{cap:caracteristicas}) & \var{python -m brugada.experiments.run\_features} \\
Reducción de dimensionalidad (\cref{cap:visualizacion}) & \var{python -m brugada.experiments.run\_dimred} \\
Modelo de referencia (\cref{cap:modelo}) & \var{python -m brugada.experiments.run\_baseline} \\
Figuras de este documento & \var{python ../memoria/make\_figures.py} \\
\bottomrule
\end{tabularx}
\end{center}

Cada experimento guarda en \var{results/} sus métricas por \emph{fold}, el resumen y las predicciones fuera de muestra de cada paciente en cada repetición, de modo que cualquier tabla o figura puede regenerarse sin volver a entrenar.
